\vspace{-0.5mm}
\section{Discussion}
\vspace{-0.5mm}
\label{sec:discuss}
In~\cref{sec:exp} and~\cref{sec:deep_analysis}, we identified key limitations of RLVR in enhancing LLM reasoning capabilities. In this section, we explore possible underlying factors that may explain why RLVR remains bounded by the reasoning capacity of the base model.

\textbf{Discussion 1: Key Differences Between Traditional RL and RLVR for LLMs are Vast Action Space and Pretrained Priors.}
Traditional RL such as AlphaGo Zero and the DQN series~\cite{silver2017alphagozero,mnih2015dqn,yue2023value} can continuously improve the performance of a policy in environments like Go and Atari games \textit{without an explicit upper bound}. There are two key differences between traditional RL and RLVR for LLMs.
First, the action space in language models is exponentially larger than that of Go or Atari games~\cite{rl4nlp}. 
RL algorithms were not originally designed to handle such a vast action space, which makes it nearly impossible to explore the reward signal effectively if training starts from scratch. 
Therefore, the second distinction is that RLVR for LLMs starts with a pretrained base model with useful prior, whereas traditional RL in Atari and GO games often begins from scratch. 
This pretrained prior guides the LLM in generating reasonable responses, making the exploration process significantly easier, and the policy can receive positive reward feedback.

\textbf{Discussion 2: Priors as a Double-Edged Sword in This Vast Action Space.}
Since the sampling of responses is guided by the pretrained prior, the policy may struggle to explore new reasoning patterns beyond what the prior already provides. 
Specifically, in such a complex and highly combinatorial space, most responses generated by \textit{naive token-level sampling exploration} are constrained by the base model's prior. 
Any sample deviating from the prior is highly likely to produce invalid or non-sensical outputs, leading to negative \textit{outcome reward}.
As discussed in~\cref{sec:prelim}, policy gradient algorithms aim to maximize the log-likelihood of responses within the prior that receive positive rewards, while minimizing the likelihood of responses outside the prior that receive negative rewards. 
As a result, the trained policy tends to produce responses already present in the prior, constraining its reasoning ability within the boundaries of the base model.
From this perspective, training RL models from a distilled model may temporarily provide a beneficial solution, as distillation helps inject a better prior. 

\textbf{Possible Future Work.} As discussed above, inefficient exploration mechanisms in a vast action space and the reliance on binary outcome rewards may be the root causes of the limitations observed in current RLVR settings. To fundamentally address these challenges, several directions may be worth exploring:

\vspace{-1.0mm}
\begin{itemize}
\item \textbf{Efficient exploration strategies in high-level abstraction.} High-level exploration mechanisms such as AlphaEvolve~\cite{novikov2025alphaevolve}, which perform self-evolution in a program-level abstraction space, may be crucial for navigating the vast action space. Such strategies could facilitate the discovery of out-of-prior reasoning patterns and previously unseen knowledge structures.

\item \textbf{Data scale via curriculum.} A curriculum can begin by training on easier subproblems, allowing the model to improve sampling efficiency and acquire essential meta-skills. By increasing success rates on simpler tasks before tackling harder ones, such a curriculum may hierarchically reduce the exploration space and lift performance from nearly zero to non-zero on challenging parent tasks, thereby enabling RLVR to obtain meaningful rewards~\cite{zhang2025stephint,li2025questa}. Although traces of such hierarchical relationships occasionally appear in current RLVR training data, and their effects have been observed in recent work~\cite{chen2025acereason}, realizing their full potential will require a more deliberate, large-scale data-RL iteration pipeline that ensures sufficient coverage of meta-skills as well as appropriate relationships between easy and hard problems.

\item \textbf{Process reward and fine-grained credit assignment.} Compared to purely binary outcome rewards, incorporating intermediate signals to guide the reasoning trajectory may significantly improve exploration efficiency and steer exploration toward more promising solution paths.

\item \textbf{Agentic RL.} Current RLVR reasoning are limited to single-turn response, whereas iterative refinement based on feedback is crucial for IMO-level reasoning~\cite{huang2025winning}. It also lacks the ability to actively collect new information by using search tools or conducting experiments. A multi-turn agentic RL paradigm, featuring richer interactions with environment feedback, could allow models to generate novel experiences and learn from them. This emerging agent framework has been described as the beginning of an "era of experience"~\cite{silver2025welcome}.


\end{itemize}

\vspace{-0.5mm}
\section{Related Work}
\vspace{-0.5mm}

We summarize key related works on the analysis of RLVR here and provide a more comprehensive discussion in Appendix~\ref{appendix:more-related}.
While recent RLVR methods have achieved impressive empirical results~\cite{guo2025deepseek-r1, lambert2024tulu3}, their fundamental impact on reasoning remains underexplored. 
Several studies~\cite{liu2025there, zhao2025echo, shah2025rethinking} suggest that reflective behaviors in RLVR models originate from the base models rather than being learned through reinforcement learning. Dang~\etal~\cite{dang2025assessing} observed a decline in pass@$k$ performance post-RLVR training, but their analysis was limited in scope. 
More importantly, they did not explore the relationship between the base model and the RL model.
Deepseek-Math~\cite{shao2024deepseekmath} also observed similar trends, but their study was limited to a single instruction-tuned model and two math benchmarks.
In contrast, our work systematically investigates a wide range of models, tasks, and RL algorithms to accurately assess the effects of current RLVR methods and models. We further provide in-depth analyses, including accuracy distributions, reasoning coverage, perplexity trends, and comparison against distilled models, offering a comprehensive understanding of RLVR's capabilities and limitations.
